% TOP_PROBLEM: {"id": 4400007, "title": "Pingree open problems — Boyle problem 1", "category_id": 11, "category": {"id": 11, "name": "dynamical_systems", "display_name": "Dynamical Systems", "description": "Problems about long-term behavior of deterministic systems, Hamiltonian dynamics, and periodic orbits.", "slug": "dynamical-systems", "order_index": 11, "created_at": "2026-07-31T15:26:25.671Z"}, "status": "open", "problem_number": "AMR-043-0007", "set_id": 13}
% FIRST_POSTED: 2026-10-01
% ENTRY_KIND: statement_only
% UnsolvedMath ID 4400007; source catalogue code AMR-043-0007
% !TeX program = lualatex
% Definition and sources only; this file is not an LLM solution attempt.
\documentclass[11pt,a4paper]{article}
\usepackage[margin=25mm]{geometry}
\usepackage{fontspec}
\setmainfont{lmroman10-regular.otf}[BoldFont=lmroman10-bold.otf,
 ItalicFont=lmroman10-italic.otf,BoldItalicFont=lmroman10-bolditalic.otf]
\usepackage{amsmath,amssymb,mathtools}
\usepackage{unicode-math}
\setmathfont{latinmodern-math.otf}
\AtBeginDocument{\let\setminus\smallsetminus}
\usepackage{xurl,hyperref}
\hypersetup{colorlinks=true,urlcolor=blue}
\setcounter{secnumdepth}{0}
\setlength{\parindent}{0pt}
\setlength{\parskip}{5pt}
\setlength{\emergencystretch}{3em}
\begin{document}
\section*{Pingree open problems — Boyle problem 1}\label{4400007}
{\small UnsolvedMath ID 4400007 · AMR-043-0007 · Dynamical Systems · AMR Open Problem Lists}
\subsection{Definitions and mathematical statement}
A two-sided subshift is a closed subset $X\subset A^{\mathbb Z}$, for a
finite discrete alphabet $A$, invariant under the homeomorphism
$(\sigma x)_j=x_{j+1}$ and its inverse. The topology is the product topology.
It is a shift of finite type (SFT) if a finite list of forbidden finite words
defines $X$. It is mixing if for any nonempty relatively open $U,V\subset X$
there is $N$ such that $\sigma^n(U)\cap V\ne\varnothing$ for every $n\ge N$.

Two such systems $(X,\sigma)$ and $(Y,\tau)$ are topologically orbit equivalent
if there is a homeomorphism $H:X\to Y$ satisfying
\[
 H(\{\sigma^n x:n\in\mathbb Z\})
   =\{\tau^n H(x):n\in\mathbb Z\}\qquad(x\in X).
\]
No requirement that $H\sigma=\tau H$ is imposed, and the integers relating
successive images need not form a continuous cocycle. The target is a
classification of mixing SFTs under this equivalence: give a complete set of
invariants, or necessary and sufficient conditions on their finite graph
presentations, for such an $H$ to exist.

An orbit homeomorphism preserves the cardinality of every finite orbit;
therefore equal numbers of periodic orbits of each length are necessary.
The classification must account for infinite orbits and their accumulation
as well. In particular, merely matching finite orbit counts does not by
itself provide the required homeomorphism of the compact spaces.
\subsection{Short English statement}
Classify mixing shifts defined by finitely many local rules when a homeomorphism is allowed to rearrange time along each orbit.
\subsection{Sources}
\begin{itemize}
\item[S1] ULAM.AI and the UnsolvedMath Contributors, supplied problems.json, record 4400007 (AMR-043-0007), AMR Open Problem Lists. Catalogue identity and source transcription; CC BY 4.0.
\url{https://huggingface.co/datasets/ulamai/UnsolvedMath}.
\item[S2] Open problems from the 3rd Pingree Workshop on Dynamical Systems. (2010), Mike Boyle, Problem 1. Original source indexed by AMR-043-0007.
\url{https://math.huji.ac.il/~mhochman/open-problems/pingree-open-problems.pdf}.
\end{itemize}
\end{document}
